
\appendix
\section{Finite coverage tables for \texorpdfstring{$t_4$ and $t_5$}{t4 and t5}}
\label{app:finite-tables}

This appendix records the finite data underlying Propositions~\ref{prop:t4-extended} and~\ref{prop:t5-extended}.  Because Proposition~\ref{prop:minimal-coverage-reduction} reduces the attractor property to the inclusion-minimal coverage antichain, these tables are complete finite certificates rather than merely collections of necessary conditions.

\subsection{The twelve minimal coverages of \texorpdfstring{$t_4$}{t4}}

For $t_4$, there are $100$ distinct nonempty factors and $99$ distinct factor coverages.  Exactly twelve of the latter are inclusion-minimal.  They are the coverages of the eight MUSs of $t_4$ and the four relevant MUSs of $t_3$:

\begin{center}
\begin{tabular}{ll@{\qquad}ll}
\toprule
factor & coverage & factor & coverage\\
\midrule
$bbab$ & $[2,5]$ & $baba$ & $[3,6]$\\
$abaa$ & $[4,7]$ & $aabb$ & $[6,9]$\\
$bbaa$ & $[8,11]$ & $aaba$ & $[10,13]$\\
$abab$ & $[11,14]$ & $babb$ & $[12,15]$\\
$aa$ & $[6,7]\cup[10,11]$ & $bb$ & $[2,3]\cup[8,9]\cup[14,15]$\\
$aba$ & $[4,6]\cup[11,13]$ & $bab$ & $[3,5]\cup[12,14]$\\
\bottomrule
\end{tabular}
\end{center}

Every four-position hitting set of these twelve coverages contains one point in $[2,5]$, two in $[6,11]$, and one in $[12,15]$.  Expanding the twelve intersection predicates and removing redundant implications yields the ten clauses of Proposition~\ref{prop:t4-extended}.  Direct counting by the value of $\ell$ gives $10+21+28+28=87$ solutions.

\subsection{The eighteen minimal coverages of \texorpdfstring{$t_5$}{t5}}

For $t_5$, there are $392$ distinct nonempty factors and $391$ distinct factor coverages.  Exactly eighteen are inclusion-minimal.  The eight current-generation coverages are

\begin{center}
\begin{tabular}{ll@{\qquad}ll}
\toprule
factor & coverage & factor & coverage\\
\midrule
$abaabb$ & $[4,9]$ & $aabbaa$ & $[6,11]$\\
$bbaaba$ & $[8,13]$ & $babbab$ & $[12,17]$\\
$abaaba$ & $[16,21]$ & $babbaa$ & $[20,25]$\\
$bbaabb$ & $[22,27]$ & $aabbab$ & $[24,29]$\\
\bottomrule
\end{tabular}
\end{center}

The remaining ten minimal coverages are

\begin{center}
\begin{tabular}{ll}
\toprule
factor & coverage\\
\midrule
$aaba$ & $[10,13]\cup[18,21]$\\
$abab$ & $[11,14]\cup[19,22]$\\
$babb$ & $[12,15]\cup[20,23]$\\
$bbaa$ & $[8,11]\cup[22,25]$\\
$aabb$ & $[6,9]\cup[24,27]$\\
$bbab$ & $[2,5]\cup[14,17]\cup[26,29]$\\
$baba$ & $[3,6]\cup[15,18]\cup[27,30]$\\
$abaa$ & $[4,7]\cup[16,19]\cup[28,31]$\\
$aa$ & $[6,7]\cup[10,11]\cup[18,19]\cup[24,25]\cup[30,31]$\\
$bb$ & $[2,3]\cup[8,9]\cup[14,15]\cup[22,23]\cup[26,27]$\\
\bottomrule
\end{tabular}
\end{center}

Every four-position transversal of these eighteen masks contains one point in $\{8,9\}$ and one in $\{24,25\}$.  Its central pair belongs to one of the three classes
\[
 \{12,13\}\times\{16,17\},
 \qquad
 \{(14,18),(15,19)\},
 \qquad
 \{16,17\}\times\{20,21\}.
\]
The outer choices are independent.  The three classes therefore produce $16$, $8$, and $16$ attractors, respectively, giving the total $40$.


\section{Details of the independent verification}
\label{app:independent-verification}

The independent verifier is intentionally separate from the implementation used during the derivation of the symbolic formulas.  Its construction follows the definitions directly and does not import helper routines from the original program.

For each tested order it generates the Thue--Morse word, enumerates every distinct factor, computes all overlapping occurrences, builds the corresponding coverage set, removes duplicate coverage sets, and then filters the remaining family by inclusion.  For the transversal computation it branches only on positions belonging to a currently unhit constraint; this is an exhaustive hitting-set search, not a heuristic one.

Besides the values in Table~\ref{tab:independent-computations}, the verifier performs the following independent checks:
\begin{itemize}
\item the complete $t_4$ and $t_5$ antichains in Appendix~\ref{app:finite-tables};
\item equality of the exhaustive attractor sets with the descriptions in Propositions~\ref{prop:t4-extended} and~\ref{prop:t5-extended};
\item the trimmed MUS identity and full occurrence transport through ambient order ten;
\item the centered/paired-boundary property through $t_9$;
\item every proposed stable attractor against every distinct factor for $6\le n\le9$;
\item the minimal-coverage theorem for $5\le n\le9$;
\item the short critical-factor lists and critical descent through order nine;
\item the depth-two counterexample and the five-position limitation of the stable MUS certificate.
\end{itemize}

No random sampling is used in any of these tests.


\section{Indexing conventions in the cited Thue--Morse literature}
\label{app:indexing}

Two indexing conventions appear in the literature used in this paper.  Kutsukake et al.~\cite{Kutsukake2020} use the same convention as we do:
\[
 t_0=a,
 \qquad
 |t_n|=2^n.
\]
Kishi et al.~\cite{Kishi2026} write
$
 T_i=\mu^{i-1}(a),
$
so that
$
 T_i=t_{i-1}.
$
Consequently their formula for the MUSs of $T_i$ translates into our Theorem~\ref{thm:kishi-mus} after the substitution $m=i-1$.  The shift is also consistent with the lengths $|T_i|=2^{i-1}$ and $|t_n|=2^n$.
